\section{Conclusiones}

El mínimo real creció 13,5\% entre 2021 y 2025, frente a 1,2\% de la productividad por hora. El aumento nominal de 23\% de 2026 amplía la ganancia salarial, pendiente de las cifras de cierre. La mejora del ingreso de quienes reciben el mínimo debe evaluarse junto con las oportunidades de contratación y sus obligaciones fiscales.

El saldo de empleo no basta para esa evaluación. La evidencia de Antioquia muestra que el ajuste puede ocurrir mediante empleos que dejan de crearse. También deben observarse el acceso al trabajo formal y la movilidad salarial. El reescalamiento nacional presentado ilustra la magnitud que resultaría de trasladar proporcionalmente el efecto regional. Sus supuestos no están verificados para el país y no permite atribuirle una pérdida observada de puestos ni medir la compresión salarial de toda Colombia.

Proponemos que el ajuste de 2027 compense la inflación de 2026 sin añadir puntos porcentuales. Con los pronósticos considerados, corresponde a 6,7\%--6,9\%, con 7\% como límite provisional sujeto a revisión si la inflación supera esa referencia. Moderar el ajuste daría margen para absorber las ganancias anteriores. Reducir la relación entre el mínimo real y la productividad requiere que esta última crezca más rápido que el primero.

La decisión de diciembre debe actualizar los pronósticos y valorar los demás criterios legales. También debe reconocer que el costo del mínimo incluye pensiones, coberturas y contratación pública, además de nómina. La protección del ingreso debe considerar también las oportunidades de quienes buscan empleo formal.

\fuente{Las definiciones y los cálculos se documentan en la \href{https://github.com/oliverpardo1979/Informe-CJC-SM/blob/main/nota_tecnica.pdf}{nota técnica de fuentes y medición}, que acompaña este informe.}
